\section{検証（Validation \& Verification）}
\label{sec:validation}

少数条件を高忠実度で解く密結合方式では，各本番実行の信頼性を事前にどれだけ担保できるかが重要になる。本章では，段階的な検証計画とその結果を要約する。

\subsection{V\&V-1：コード基礎検証}

層流Poiseuille流，Taylor-Green渦など解析解を持つベンチマークで，LBM実装の収束次数を確認した。減衰Taylor-Green渦（3次元格子にz方向一様で埋め込んだ2次元渦）で，実効動粘性が\eref{viscosity}と一致すること（運動エネルギー減衰率から逆算して誤差2\%以内），速度場が解析解と2\%以内で一致すること（相対L2誤差，\(n=32\)，200ステップ），空間収束次数が設計どおり2次精度（\(n=16/32/64\)で2.03/1.94，\(\tau=0.8\)）であることを確認した。

\(\tau\)を\(1/2\)近傍（例：0.55）に取ると，BGKの高次誤差項が2次項と同程度になり測定次数が1.3--1.8へ乱れることも確認しており，収束次数の検証は\(\tau=0.8\)程度の条件で行う必要がある。

\subsection{V\&V-2：乱流モデル検証（クリープ流極限）}

周期立方配列中の球（\(a/L=1/6\)）に体積力を与え，定常状態の抗力をHasimotoの解析解と比較した。レイノルズ数を一定（\(\approx0.03\)）に保ち解像度のみを変えたところ，補間バウンスバック＋厳密な\(\delta\)評価では，実測抗力/Hasimoto予測の比が半径4--8格子で0.958--0.974と，解像度とともに単調に1へ近づくことを確認した。定常状態での運動量収支（壁への力／体積力の総和）は全ケースで0.9999以上であり，力の算出そのものが正しいことも確認済みである。

\subsection{V\&V-2：乱流モデル検証（高レイノルズ数側）}
\label{sec:validation-highre}

一様流中の滑面球についてClift-Gauvin相関との比を測定した。当初，周期境界のみの構成では，低レイノルズ数域で比が0.5未満と大きく外れる現象が見られた。詳細な切り分けの結果，これは離散化誤差ではなく\emph{周期境界による後流の巻き込み}が原因であることが判明した：周期箱では球に届く流れが箱平均速度の31\%程度しかなく，抗力係数の基準速度が大きく過大評価されていた。

\subsubsection{流入・流出境界の導入とその効果}

\secref{open-boundary}の流入・流出境界を導入し，同じ条件（\(Re=100\)，N/D=20，横幅4D）で再測定したところ，球に届く流れは入口指定値の96.4\%まで回復し，Clift-Gauvin比は0.47から1.12（相関自身の不確かさの範囲内）まで改善した。幅を無限大へ線形外挿すると，比は\textbf{1.012}に収束する。この結果により，\(Re=100\)におけるV\&V-2は合格と判断した。

\subsubsection{壁面条件の検証}

流入・流出境界のもとで，領域幅・バウンスバック規則・緩和時間をそれぞれ独立に変えて壁面条件の妥当性を検証した。

\begin{table}[H]
	\centering
	\caption{バウンスバック規則の比較（幅4D，Clift-Gauvin比）。\(\tau\)・Ma・Re・N/Dはすべて同一で，壁位置のみが異なる。}
	\label{tab:bounceback-rule}
	\begin{tabular}{lcc}
		\Hline{1.5pt}
		規則 & N/D=20 & N/D=40 \\
		\hline
		補間バウンスバック & 1.12 & 1.09 \\
		ハーフウェイ & 1.16 & 1.11 \\
		\hline
		規則差 & 3.5\% & \textbf{1.8\%} \\
	\end{tabular}
\end{table}

規則差は\(\Delta x\)を半分にすると半減しており，壁位置の1次の離散化誤差として設計どおり収束している。この結果に基づき，補間バウンスバック規則を採用した。

中心モーメント演算子では，3次以上のキュムラントの緩和率が2次モーメントのそれと独立に選べるため，TRTのmagic parameter \(\Lambda=(1/\omega^+-1/2)(1/\omega^--1/2)\)を自由に設定できる（\secref{cumulant}）。ポアズイユ流の速度分布から実効壁位置を放物線フィットで取り出して測定したところ，既定値 \(\omega^-=1\) が実質的に最適であり，\(\Lambda\)をTRTの理論値\(3/16\)（壁を格子の中間に置くとされる値）に強制すると，むしろ壁位置は数倍悪化することを確認した。この結果は，壁面条件が\(Re=100\)で観測される緩和時間依存性（後述）の原因ではないことを示す独立した証拠になっている。

\subsubsection{緩和時間依存性の切り分け}

流入・流出境界のもとでも，緩和時間 \(\tau\) を振ると抗力比が7\%程度動く現象が残った。これに対し，圧縮性（\(Ma^2\)），壁面条件（上述），解像度で消える打ち切り誤差，数値粘性の下限，単精度の丸め誤差という5つの仮説を，いずれも「動かせば動くはずのものを動かして，動かなかった」形で個別に検証し，すべて棄却した。原因は未特定のまま残っているが，Clift-Gauvin相関自身の不確かさとV\&V-5の許容誤差（後述，変化量20\%）を踏まえ，この6\%程度の緩和時間依存性はコード全体の不確かさとして申告した上で本番実行に進む方針とした。なお本番の\(\tau-1/2\)は，この検証で振った範囲よりもさらに一桁以上\(1/2\)に近い領域にある。

\subsection{V\&V-2：局所グリッドリファインメントの直接検証}

外側領域・ボール・タイムステップ数を揃え，一様な細格子とリファインド格子で力を比較した。

\begin{table}[H]
	\centering
	\caption{一様細格子とリファインド格子の抗力・揚力の比較。}
	\label{tab:refine-drag}
	\begin{tabular}{lrr}
		\Hline{1.5pt}
		構成 & 抗力 & 揚力 \\
		\hline
		一様細格子（基準） & \(-0.5975\) & 0.01527 \\
		リファインド，パッチ3.5D & \(-0.5905\)（\(-1.2\%\)） & 0.01467（\(-4.0\%\)） \\
		リファインド，パッチ2.5D & \(-0.6240\)（\(+4.4\%\)） & 0.01633（\(+6.9\%\)） \\
		\hline
	\end{tabular}
\end{table}

パッチが十分大きければ（本プロジェクトでは3直径以上），リファインド格子は一様細格子を数\%以内で再現する。狭いパッチが外れるのは，界面がボール表面から近い位置に来て，粗レベルが表現できない近傍後流をそこで課すことになるためである。

\subsection{V\&V-3：回転球検証}

滑面回転球のマグヌス係数について，\(S<0.4\)の範囲で\(C_L\approx S\)の線形関係を再現することを確認する（実施予定）。CPU参照実装では，無回転で横力・トルクが機械精度でゼロになること，マグヌス力の向きが\(\vec\omega\times\hat V\)と一致しスピン反転に対して反対称であること，流体が自転を引き戻す向きのトルクが出ることを確認済みである。

\subsection{V\&V-4：縫い目球検証}

縫い目付き球体を固定・一様流中に置いた場合の\(C_D\)，\(C_L\)を実測データ\cite{kensrud2010thesis,nathan2008}と比較し，縫い目向き依存性・非マグヌス力の再現性を確認する（実施予定）。密結合ループの助走・スピンアップ段階を流用できる。

\subsection{V\&V-5：軌道検証}
\label{sec:vv5}

既知の実測ピッチトラッキングデータと，密結合シミュレーションの軌道出力を比較する。検証ケースとして，WBC2023決勝，大谷翔平の対トラウト最終球（2023年3月21日，9回裏2死フルカウント）を採用した\cite{cbssports2023,si2023}。

\begin{table}[H]
	\centering
	\caption{V\&V-5検証ケースの実測値。}
	\label{tab:vv5-case}
	\begin{tabular}{ll}
		\Hline{1.5pt}
		項目 & 実測値 \\
		\hline
		球種（Statcast分類） & スウィーパー \\
		球速 & 87.2\,mph（\(\approx38.99\,\text{m/s}\)） \\
		スピン量 & 2708\,rpm \\
		横変化量 & 17\,in（\(\approx0.432\,\text{m}\)） \\
		縦変化量 & 32\,in（\(\approx0.813\,\text{m}\)） \\
		\hline
	\end{tabular}
\end{table}

球速・回転数はシミュレーションの初期条件として直接使用し，横変化量・縦変化量をシミュレーション軌道から同一定義で算出した値と比較することで合格基準とする。ただし，スピン軸の正確な向き（チルト角・ヨー角）とリリースポイント座標は本稿執筆時点で未確認であり，別途確認またはキャリブレーションで補う必要がある。

\paragraph{「変化量」の定義について}
「変化量」は単一の軌道の性質ではなく2本の軌道の差であり，どちらを基準に取るかで値が2倍近く変わる。この差は許容誤差（暫定20\%）よりはるかに大きいため，比較の前に定義を確定させる必要がある。実装では3つの定義（pfx：本塁手前40\,ft地点で分岐しマグヌス力のみ切った軌道，induced：リリース時点から分岐した同様の軌道，total：無力の直進基準）をすべて算出し，報告値がどの定義と対応するかを明示的に比較できるようにしている。
